\section*{Introduction}
\phantomsection
\addcontentsline{toc}{section}{Introduction}
\label{v:presentacion}
Ce travail établit la dérivation des distributions de
Bose--Einstein et de Fermi--Dirac à partir des complétions quantiques
d’états avec trajectoire et mémoire, et démontre leur convergence vers le régime
dilué de Maxwell--Boltzmann. La même statistique bosonique, réalisée
sur les modes transversaux du champ électromagnétique, détermine la
loi spectrale de Planck, le flux de Stefan--Boltzmann et les
déplacements de Wien. Le problème central est de relier la structure
de l’état, ses occupations admissibles et ses observables thermiques au moyen
d’opérateurs et de mesures explicites. Ainsi, les lois d’occupation et de
rayonnement s’obtiennent en des étapes reliées d’une construction, avec les
conditions d’équilibre et de représentation qui correspondent à chacune.

La base est l’holographie modulaire triadique (HMT), avec sa réalisation
dimensionnelle en mécanique dimensionnelle. La construction des chemins
\emph{All Possible Paths} (APP) est organisée sur le support
\((\mathbb Z/9\mathbb Z)^2\). Les translations par
\(\pm(1,0)\) et \(\pm(0,1)\) définissent son graphe de Cayley additif ;
les cellules carrées fournissent une structure toroïdale bidimensionnelle.
Sur ce même support, on évalue la somme et le produit, en conservant
les résidus et les quotients. La signature ternaire orientée, TRIT, prend ses valeurs
dans \(\{-1,0,+1\}\) et distingue l’orientation et le régime. Son relèvement
entier conserve conjointement le résidu et la retenue. Le noyau
topologique de phase, \emph{Topological Phase Kernel}
(TPK), compose la sélection de phase, le transport, la mise à jour et le
prolongement de l’état. Son registre retient la trajectoire et la mémoire :
le retour de phase ne réinitialise pas l’état.

Les conditions initiales et le raffinement compatible produisent les
sections généalogiques utilisées dans la réalisation linéaire. La
génération régionale, l’arithmétique des histoires, le registre à douze phases,
la détermination d’alpha et l’incidence exceptionnelle fixent les
antécédents des coefficients et de leurs représentations. Sur ce
domaine, on construit l’espace des états, l’évolution unitaire et
l’occupation multiparticulaire. La compatibilité avec les restrictions de
profondeur précède son transport entre réalisations ; dans l’instance
finie, le calendrier fournit un Hamiltonien, un propagateur et une
fonctionnelle d’action explicites.

Le pas décisif vers la statistique est le caractère compositionnel de
parité. Le secteur pair est complété au moyen de puissances symétriques, avec
des occupations entières non négatives ; le secteur impair, au moyen de puissances
extérieures, dont les opérateurs de création sont nilpotents et dont les
opérateurs d’occupation sont des projecteurs. On obtient ainsi
l’exclusion algébrique de Pauli. L’identification relativiste de cette
parité au spin conserve, en outre, la compatibilité de la
représentation des rotations et les conditions de localité. Les multiplicités déterminent les dénombrements microcanoniques, et la
factorisation des traces de Gibbs produit les distributions de
Bose--Einstein et de Fermi--Dirac. Dans le régime thermodynamique, la variation
de l’entropie combinatoire sous contraintes d’énergie et de population
fournit une seconde caractérisation de l’équilibre. Pour le secteur bosonique, on maintient la
condition sur le potentiel chimique qui garantit la convergence.
Les bornes du régime dilué quantifient la limite commune de
Maxwell--Boltzmann.

La réalisation radiative applique l’occupation bosonique à une cavité
périodique tridimensionnelle avec deux polarisations transversales, une dispersion
linéaire et un potentiel chimique nul. Le mode constant est fixé comme donnée de
fond avant de former l’état des excitations. La densité de modes et
l’énergie par excitation produisent le spectre de Planck. La convergence
à partir des sommes discrètes est démontrée en contrôlant uniformément le
voisinage de la fréquence nulle et la queue des hautes fréquences. Les moments
spectraux donnent les densités de nombre, d’énergie et d’entropie ;
l’intégration angulaire donne le flux de Stefan--Boltzmann. Le changement entre
fréquence et longueur d’onde transforme aussi la mesure au moyen de son
jacobien, de sorte que les deux maxima de Wien satisfont des équations
distinctes. Chaque constante de déplacement est caractérisée par le
maximum de la densité correspondante.

La relation entre action, calendrier et température détermine les
normalisations dimensionnelles de cette réalisation. La transduction
thermique distingue l’énergie de décision, l’information sans dimension et
la paire Boltzmann--température. Sa composition avec la vitesse
constitutive relie les échelles temporelle, énergétique et spatiale du
rayonnement. Les hypothèses de continuité, de domaine du générateur,
de contractivité et de convergence restent associées aux opérations
qui les requièrent.

Entre la transduction et les complétions, la chaîne qui fixe le porteur
se compose du renouvellement tritique, de la double graduation, de la mémoire énergétique positive,
du spectre de Williamson, du raffinement par trois, de la section marquée, des courants de
bord et de l’équivalence thermique centrée. Cette chaîne préserve
l’origine qui disparaît d’une densité normalisée et la transporte jusqu’à l’aire,
l’énergie libre, la gravitation et le rayonnement. C’est pourquoi l’évaluation complète n’est pas
présentée comme une coïncidence décimale, mais comme la sortie d’opérateurs, de fibres et
d’applications typés.

\Needspace{13\baselineskip}
L’interprétation entropique distingue les configurations d’occupation,
la mesure sur les histoires et le spectre d’un opérateur densité.
La fréquence chronologique d’un calendrier n’est pas la mesure stationnaire
de son enveloppe d’histoires. Cette dernière est caractérisée
variationnellement au moyen d’une identité de divergence relative.
Dans le secteur fermionique quasi libre conservant le nombre et invariant
sous les transformations de jauge, le corrélateur à un corps \(0<C<I\)
et son générateur modulaire \(K_1\) se déterminent réciproquement en dimension
finie. Le transport, la
réduction et le conditionnement précisent quelle information demeure
observable. La récupération du registre et l’identification orbitale des
opérateurs s’articulent avec un bilan bilatéral qui conserve
conjointement l’état et la mémoire
(section~\ref{sec:acoplamiento-recuperacion}).

L’unité surfacique et l’intrication complètent cette relation entre
structure et information. L’identité modulaire finie exprime les
variations d’entropie au moyen des aires sectorielles et d’un déficit
d’entropie relative non négatif qui est soustrait ; avec \(\Delta_{\rm mod}\ne0\) et
une normalisation connue, l’aire et le Hamiltonien modulaire permettent de
récupérer les occupations des deux secteurs. Le lien entre mémoire, occupation et
rayonnement est donc établi en conservant les applications qui permettent de
passer d’un objet au suivant, en plus de leurs valeurs scalaires.
L’appendice~\ref{v:integracion} spécifie les domaines des
réalisations et les conditions de ces compositions.

La réalisation finie sur \(\Lambda=(\mathbb Z/27\mathbb Z)^2\)
établit des bornes exactes de support et d’entropie, avec saturation par
des sous-groupes et des annulateurs. Dans le domaine dodécaphasique, \(\nu_{120}\) et
\(\nu_{270}\) détectent des corrélations d’ordre que l’histogramme tritique
élimine. Sur le domaine premier, \(R_{\rm vac}\) compose le mot
sturmien engendré, la réalisation interne des irréductibles et l’horloge
\(\log p\). Son identité modale est comprise au moyen de sommes symétriques ou de
moyennes de Fejér, et ses tableaux constituent des témoins numériques finis.

\clearpage
\subsection*{Relations entre les constructions et leurs résultats}
Le tableau suivant identifie la contribution de chaque bloc à l’argument.
L’ordre d’exposition conserve les antécédents communs, puis distingue
les opérations spécifiques : représentation, occupation, conditionnement et
lecture radiative. La représentation graduée détermine les occupations ; la réalisation
spatiale et la dispersion déterminent leur application radiative.

\begin{table}[H]
\centering
\small
\renewcommand{\arraystretch}{1.2}
\begin{tabularx}{\textwidth}{@{}>{\raggedright\arraybackslash}p{.29\textwidth}>{\raggedright\arraybackslash}X@{}}
\toprule
Construction et localisation & Construction et fonction dans l’argument \\
\midrule
Noyau et prolongements, §§\ref{sec:nucleo}--\ref{exc:seccion}
& États admissibles, chemins orientés, mémoire et raffinement ; sorties
régionales, registre dodécaphasique et incidence. Ces antécédents conservent
l’origine des coefficients et des représentations ultérieures. \\
Réalisation quantique et action,
§§\ref{v:rec:sec:postulados-cuanticos-genealogicos-rev2}--\ref{sec:hmtf2-cristal-temporal}
& Espace des états, évolution unitaire et action du calendrier. La
réalisation fixe les opérateurs qui interviennent ensuite dans les énergies
et dans les états d’équilibre. \\
Mémoire et température,
§\ref{v:ant:sec:ii-kb-aut-desarrollo}
& Fibres de lecture, réversibilité et paire Boltzmann--température.
On distingue l’information sans dimension, l’échelle énergétique et le système de coordonnées thermiques. \\
Capacité spectrale et pont marqué,
§\ref{v:sec:capacidad-espectral-termica}
& Arbre microscopique pondéré, porteur gradué, courants de bord,
échelles de capacité et de Williamson, et transducteur thermométrique variationnel. \\
Parité, complétions de Fock et distributions d’équilibre,
§\ref{v:sec:completaciones-algebraicas} et
§§\ref{v:rec:fock_gibbs:section:01}--\ref{v:rec:ocupacion:section:03} ;
jalons \ref{v:rec:completaciones:statement:02},
\ref{v:rec:fock_gibbs:statement:01},
\ref{v:rec:thm:factorizacion-gran-canonica-rev3} et
\ref{v:rec:thm:limite-clasico-fd-be-rev2}
& Parité, Fock, transport par raffinement et état de Gibbs.
Les multiplicités donnent les dénombrements ; les traces produisent les deux
distributions d’occupation, dont la limite diluée est quantifiée, §§\ref{v:rec:ocupacion:section:02}--\ref{v:rec:ocupacion:section:03}. \\
Histoires, aire et\newline intrication
& L’incidence surfacique--volumétrique détermine la mesure des histoires ;
la réduction des états et l’opérateur d’aire soutiennent la lecture
modulaire de l’entropie. Voir le théorème~\ref{thm:incidencia-hojas-anexo}
et la section~\ref{v:ant:sec:area}. \\
Incertitude de Fourier, mémoire ordonnée et vacances premières,
§\ref{v:n19:sec:incertidumbre-fourier} et
§\ref{sec:f051-prime-vacancy-observer}
& La réalisation finie établit les bornes de support et d’entropie ; les
lecteurs dodécaphasiques distinguent l’ordre et les marginales ; l’observateur premier
des vacances résout la corrélation au moyen de l’identité modale
symétrique/Fejér, avec un témoin historique jusqu’à \(10^8\), une reproduction
indépendante jusqu’à \(10^6\) et une borne critique conditionnée par
l’uniformité modale. \\
Occupation bosonique et lois du rayonnement thermique,
§\ref{v:sec:radiacion-termica} ; propositions
\ref{v:prop:ley-espectral}, \ref{v:prop:stefan} et \ref{v:prop:wien} ;
limite du lemme~\ref{v:lem:limite-termodinamico}
& Les modes transversaux, la vitesse constitutive et la statistique
bosonique déterminent les sommes thermiques. Les estimations à l’origine et sur la queue
justifient la limite ; leurs moments fournissent les densités, le flux et les
maxima spectraux. \\
\bottomrule
\end{tabularx}
\caption{Constructions précédentes et utilisation ultérieure dans l’argument
statistique et radiatif.}
\label{v:tab:mapa-demostrativo}
\end{table}

La présence d’antécédents partagés répond à cette dépendance effective.
Chaque bloc conserve l’information dont l’opération suivante a besoin ;
la coïncidence d’une dimension ou d’une valeur scalaire s’accompagne de l’application
qui permet de la réutiliser. Ainsi, la statistique reçoit une représentation
graduée et un Hamiltonien, tandis que le rayonnement reçoit en outre une
réalisation spatiale et sa dispersion.
